\def\stat{kor+sor}





\def\tit{СЕРВИС-ОРИЕНТИРОВАННЫЙ ИНТЕРФЕЙС ДЛЯ~ДОСТУПА 


К~НАУЧНЫМ ДАННЫМ В~ОБЛАСТИ ИССЛЕДОВАНИЯ 


И~ОПЕРАТИВНОГО МОНИТОРИНГА СОСТОЯНИЯ ВУЛКАНОВ 


КАМЧАТКИ И~СЕВЕРНЫХ КУРИЛ$^*$}





\def\titkol{Сервис-ориентированный интерфейс для~доступа 


к~научным данным} % в~области исследования  и~оперативного мониторинга состояния вулканов  Камчатки и~Северных Курил}








\def\aut{С.\,П.~Королёв$^1$, И.\,М.~Романова$^2$, С.\,И.~Мальковский$^3$, 


А.\,А.~Сорокин$^4$}





\def\autkol{С.\,П.~Королёв, И.\,М.~Романова, С.\,И.~Мальковский, 


А.\,А.~Сорокин}





\titel{\tit}{\aut}{\autkol}{\titkol}





\index{Королёв С.\,П.}


\index{Романова И.\,М.}


\index{Мальковский С.\,И.}


\index{Сорокин А.\,А.}


\index{Korolev S.\,P.}


\index{Romanova I.\,M.}


\index{Malkovsky S.\,I.} 


\index{Sorokin A.\,A.}





{\renewcommand{\thefootnote}{\fnsymbol{footnote}}


\footnotetext[1]


{Работа выполнена при поддержке гранта РФФИ №\,16-37-00026~мол\_а 


и~Комплексной программы 


фундаментальных научных исследований ДВО РАН <<Дальний Восток>> (проект 18-5-091).


 Для обеспечения 


функционирования научных информационных сервисов были использованы ресурсы 


ЦКП <<Центр данных  ДВО РАН>>.}}





\renewcommand{\thefootnote}{\arabic{footnote}}


\footnotetext[1]{Вычислительный центр Дальневосточного отделения Российской академии наук, 


\mbox{serejk@febras.net}}


\footnotetext[2]{Институт вулканологии и~сейсмологии Дальневосточного 


отделения Российской академии наук,  \mbox{roman@kscnet.ru}}


\footnotetext[3]{Вычислительный центр Дальневосточного отделения Российской академии наук, 


\mbox{sergey.malkovsky@gmail.com}}


\footnotetext[4]{Вычислительный центр Дальневосточного отделения Российской академии наук, 


\mbox{alsor@febras.net}}








 \label{st\stat}





                  \thispagestyle{headings}


                  


  \vspace*{-8pt}


     





\Abst{На территории Дальнего Востока России располагаются десятки активных вулканов, 


требующих непрерывного внимания ученых для анализа и~контроля их состояния. 


В~качестве вспомогательных средств в~этой работе выступают информационные системы


(ИС), 


обеспечивающие решение различных научных задач. Однако разрозненность 


ИС и~ограниченный доступ к~информации существенно ограничивают 


возможность проведения комплексных исследований, что чревато в~итоге 


катастрофическими последствиями для населения и~народного хозяйства.


      Представлено описание разработанного сер\-вис-ори\-ен\-ти\-ро\-ван\-но\-го 


программного интерфейса, реализующего взаимодействие между основными 


существующими ИС для взаимного использования накопленных 


наборов научных данных и~средств их обработки при проведении исследований 


и~оперативного мониторинга со\-сто\-яния вулканов Камчатки и~Северных Курил.}





\KW{вулканы Камчатки; мониторинг вулканов; веб-сервисы; REST-сервис; интеграция 


данных; базы данных}





\DOI{10.14357\!/\!08696527180207}








\vskip 10pt plus 9pt minus 6pt








      \thispagestyle{headings}


      


      \vspace*{-10pt} 





\section{Введение}





     В настоящее время исследование окружающей среды является одним из 


тех направлений на\-уч\-но-тех\-ни\-че\-ской деятельности, которые формируют 


объективные потребности в~создании новых методов работы с~научными 


данными и~разработки на их основе современных информационных технологий 


и~систем. В~первую очередь, это связано с~активным развитием средств 


инструментальных наблюдений (наземных сетей наблюдения, систем 


дистанционного зондирования Земли и~т.\,п.)\ и~непрерывным 


совершенствованием измерительной аппаратуры. Разнообразие 


фундаментальных и~прикладных научных задач порождает создание 


разрозненных, отличных по масштабу и~функциональным возможностям 


специализированных ИС. В~то же время проведение 


междисциплинарных исследований требует совместного анализа всех 


доступных массивов неоднородной информации, сформированных 


с~применением различных программных средств и~технологий. Возникает 


потребность в~обмене данными и~организации доступа к~ресурсам внешних 


гетерогенных источников тематической информации.





Перечисленные выше утверждения справедливы для Дальнего Востока России, 


где расположено~70~действующих вулканов, требующих ежедневного 


мониторинга со стороны ученых: 30~на Камчатке и~40~на Курилах~[1]. Для их 


исследования и~систематизации накопленных данных созданы 


ИС, решающие широкий круг научных задач. Однако 


разрозненность справочной и~инструментальной информации, ограниченный 


доступ к~результатам работ существенно ограничивают возможность 


проведения комплексных исследований, а также увеличивают время на анализ 


возникающих природных событий и, как следствие, увеличивают время 


реакции на них, что чревато катастрофическими последствиями для населения 


и~народного хозяйства.


     


     В~статье представлено описание разработанного  


сер\-вис-ори\-ен\-ти\-ро\-ван\-но\-го программного интерфейса, реализующего 


взаимодействие между основными существующими 


ИС для взаимного использования накопленных наборов научных 


данных и~средств их обработки при проведении исследований и~оперативного 


мониторинга состояния вулканов Камчатки и~Северных Курил.





\section{Постановка задачи }





     Для проведения ретроспективных исследований и~оперативного 


мониторинга вулканической активности на Дальнем Востоке России в~разное 


время были созданы различные специализированные ИС. 


Если в~начале 2000-х~гг.\ они представляли собой ин\-тер\-нет-сай\-ты со 


справочными материалами и~результатами различных экспедиций 


и~исследований, то с~развитием научной телекоммуникационной~[2] 


и~инструментальной (наземной, спутниковой) инфраструктуры наблюдений 


стали появляться интерактивные сервисы для сбора, обработки и~анализа 


различной измерительной информации. Среди них можно выделить несколько 


ИС, обладающих уникальным набором данных и~средств 


их обработки (табл.~1). 


     











Информационные системы VOKKIA и~KVERT являются источниками базовой метаинформации по 


исследуемым объектам и~возникающим событиям, а~ИС <<Сигнал>> 


и~VolSatView агрегируют на своей платформе архивы различных 


инструментальных данных, а также средства для их обработки. 











     Перечисленные ИС по своим функциональным возможностям 


и~назначению дополняют друга, одновременно с~этим возникает задача обмена 


результатами исследований и~доступа к~первичной справочной информации.


\begin{table}\small %tabl1


\begin{center}


\Caption{Информационные системы  для исследования и~оперативного мониторинга 


вулканов Камчатки и~Северных Курил}


\vspace*{2ex}








\begin{tabular}{|p{40mm}|p{82mm}|}


\hline


\multicolumn{1}{|c|}{Название ИС}&\multicolumn{1}{c|}{Основное назначение}\\


\hline


1. VOKKIA (Volcanoes of  


Kurile-Kamchatka Island Arc~--- Вулканы 


Ку\-ри\-ло-Кам\-чат\-ской островной дуги)~\cite{3-kor} 


&Интеграция гетерогенных научных (описательных, количественных, пространственных) 


данных по вулканам Ку\-ри\-ло-Кам\-чат\-ской островной дуги и~их историческим 


извержениям\\


\hline


2. KVERT (Kam\-chat\-kan\linebreak Volcanic Eruption Response Team~--- Камчатская группа реагирования 


на вулканические  


извержения)~\cite{4-kor} &Автоматизация процесса подготовки сообщений KVERT 


о~текущей активности вулканов, в~том числе оперативных сообщений VONA (Volcano 


Observatory Notice for Aviation) о~пепловых шлейфах и~изменении АЦК, 


обеспечение хранения данных и~возможности их статистического анализа\\


\hline


3. <<Сигнал>>~\cite{5-kor, 6-kor} &Автоматизация операций, связанных с~работой сетей 


инструментальных наблюдений ДВО РАН, в~том чис\-ле системы видеонаблюдения за 


вулканами Камчатки и~системы компьютерного моделирования распространения пепловых 


облаков\\


\hline


4. VolSatView (Дистанционный мониторинг активности вулканов Камчатки и~Курил)~\cite{7-


kor} &Оперативный мониторинг вулканической активности на основе данных 


дистанционного зондирования Земли\\


\hline


\end{tabular}


\end{center}


%\vspace*{-12pt}


\end{table} 


Требуется организовать взаимодействие между системами, на базе которых уже 


развиты необходимые научные компетенции, сформированы архивы данных 


и~вы\-стро\-ена соответствующая вспомогательная про\-грам\-мно-ап\-па\-рат\-ная 


инфраструктура. Современным и~эффективным подходом к~решению 


указанной задачи, используемым сегодня в~ведущих мировых научных центрах 


данных (IRIS Data Management Center\footnote{{\sf http://service.iris.edu.}}, NASA 


Coordinated Data Analysis System\footnote{{\sf 


https://cdaweb.sci.gsfc.nasa.gov/WebServices/REST.}}, UCAR Research Data 


Archive\footnote{{\sf https://rda.ucar.edu/\#!apps\_api\_desc.}}), является создание  


веб-сер\-ви\-сов~\cite{8-kor, 9-kor, 10-kor}. 





Использование универсальных 


стандартов и~протокола HTTP(s) позволяет с~минимальными временн$\acute{\mbox{ы}}$ми 


затратами без внесения существенных изменений в~программный код систем 


сформировать новые виды информационных продуктов, повысить 


оперативность и~достоверность проводимых исследований.


{\looseness=1





} 


     


     На основе этих принципов и~с учетом особенностей рассматриваемых 


прикладных ИС была предложена архитектура работы 


сервиса и~разработан набор вспомогательных программных средств. 





\vspace*{-2pt}





\section{Архитектура и~принципы работы сервиса}





\vspace*{-2pt}





     Программной основой для реализации веб-сер\-ви\-са выступает ИС 


<<Сигнал>>. Из баз данных (БД) ИС VOKKIA и~KVERT в~нее поступает справочная 


информация об активных вулканах и~оперативные данные об их состоянии. 


К~оперативной информации относятся Авиационные цветовые коды (АЦК)~--- 


коды опасности вулканов для авиации, а также данные о пепловых выбросах, 


содержащиеся в~оперативных VONA-со\-об\-ще\-ни\-ях~\cite{1-kor}. Передача 


указанных данных организована посредством statement-based репликации 


соответствующих таблиц БД, где система управления БД (СУБД) ИС VOKKIA и~KVERT 


(MariaDB) является master-сер\-ве\-ром, а СУБД ИС <<Сигнал>> (MySQL)~--- 


slave-сер\-ве\-ром (рис.~1).





\begin{figure}[b] %fig1


     \vspace*{-5pt}


 \begin{center}


 \mbox{%


 \epsfxsize=126.76mm 


 \epsfbox{kor-1.eps}


 }


\end{center}


 \vspace*{-11pt}


\Caption{Общая схема работы сервиса}


\end{figure}





 Такой подход обладает следующими преимуществами:


\begin{itemize}


\item данные доступны сразу в~момент их публикации, не требуется 


периодически проверять наличие новой информации;


\item схема информационного взаимодействия серверов устойчива 


к~перерывам в~работе каналов связи благодаря заложенным в~основу 


MySQL-репликации алгоритмам;


\item программная реализация не требует больших временн$\acute{\mbox{ы}}$х затрат 


и~перерыва в~работе ИС.


\end{itemize}





     На slave-сервере ИС <<Сигнал>> исходные данные синхронизируются 


с~промежуточной БД Kvert, содержащей таблицы volcanoes и~ash. 


С~по\-мощью триггеров данные из нее передаются в~основную БД


сервиса Signal, на основе которой осуществляется обработка поступающей 


информации: изменение интервала фотосъемки вулкана в~зависимости от его 


АЦК и~постановка задания на расчет траектории распространения пеплового 


облака при извержении вулкана (рис.~2).





\begin{figure} %fig2


 \vspace*{1pt}


 \begin{center}


 \mbox{%


 \epsfxsize=125mm 


 \epsfbox{kor-2.eps}


 }


\end{center}


 \vspace*{-11pt}


\Caption{Схема данных для slave-сервера ИС <<Сигнал>>}


\end{figure}


 


     При обновлении АЦК в~таблицy Signal.volcano\_codes добавляется 


соответствующая запись. Идентификатор вулкана и~его АЦК передаются 


триггером в~хранимую процедуру, выполняющую следующую 


последовательность действий (рис.~3):


     \begin{enumerate}[(1)]


\item производится поиск всех видеокамер, осуществляющих наблюдение за 


данным вулканом;


\item для каждой видеокамеры по таблице volcano\_camera\_intervals 


определяется интервал съемки, соответствующий данному АЦК;


\item найденное значение интервала устанавливается в~качестве текущего 


для данной камеры.


\end{enumerate}











     Связи между видеокамерами и~вулканами, за которыми ведется 


наблюдение, определяются по таблице volcano\_history. Интервалы фотосъемки 


изображений с~каждой камеры могут настраиваться как независимо для каждой 


камеры (в~этом случае в~таблице volcano\_camera\_intervals поле code\_id 


принимает значение~0), так и~глобально для каждого АЦК (в~этом случае 


в~таблице volcano\_camera\_intervals поле object\_id принимает значение~0). 


Значение интервала, установленное для камеры, имеет высший приоритет по 


отношении к~глобальному значению, установленному для АЦК. 





\begin{figure} %fig3


 \vspace*{1pt}


 \begin{center}


 \mbox{%


 \epsfxsize=128mm 


 \epsfbox{kor-3.eps}


 }


\end{center}


 \vspace*{-11pt}


\Caption{Процедура БД Signal, управляющая расписанием сбора фотоснимков вулканов}


\vspace*{-6pt}


\end{figure}


     


Данные из сообщений VONA, содержащие основные характеристики каждого 


извержения (дата и~время, длительность, высота пеплового облака и~т.\,п.), 


передаются в~таблицу Kvert.ash, где на их основе для каждого события 


формируются задания на моделирование траектории движения пеплового 


облака. Задания представляют собой JSON-объ\-ек\-ты


(JavaScript Object Notation), которые с~помощью UDF-рас\-ши\-ре\-ния для СУБД 


MySQL\footnote{{\sf https://launchpad.net/gearman-mysql-udf.}} передаются на сервер очередей 


Gearman\footnote{{\sf http://gearman.org.}}, где далее обрабатываются подсистемой моделирования ИС 


<<Сигнал>>~\cite{8-kor}. Полученные результаты пуб\-ли\-ку\-ют\-ся в~сводных каталогах совместно со справочной, 


оперативной и~инструментальной информацией, доступной через пользовательский интерфейс ИС <<Сигнал>>.


     


\section{Интерфейс для доступа к~данным }





     Важной задачей веб-сервиса является организация удаленного доступа 


к~консолидированному набору данных по вулканам и~результатам работы 


специализированных подсистем ИС <<Сигнал>> для внешних пользователей, 


в~роли которых выступают сторонние ИС. Интерфейс 


взаимодействия с~ресурсами веб-сер\-ви\-са реализован в~соответствии 


с~концепцией REST~\cite{9-kor} на основе фреймворка Yii2~\cite{10-kor}. 


Работа с~ним осуществляется путем обращения по протоколу HTTP(S) к~адресу 


{\sf http://api.signal.febras.net/$\langle$url$\rangle$}, где {\sf 


$\langle$url$\rangle$}~--- относительный адрес, обозначающий вид 


запрашиваемого набора данных. На данный момент в~таком формате доступна 


информация, представленная в~табл.~2.


     


\begin{table}[b]\small %tabl2


\vspace*{-12pt}


\begin{center}


\Caption{Виды наборов данных, предоставляемых веб-сервисом}


\vspace*{2ex}





\begin{tabular}{|l|l|}


\hline


\multicolumn{1}{|c|}{URL}&\multicolumn{1}{c|}{Набор данных}\\


\hline


/volcanoes&Общий список вулканов\\


/volcanoes/\#id&Информация по вулкану с~идентификатором \#id\\


/cameras&Общий список видеокамер\\


/cameras/\#id&Информация по видеокамере с~идентификатором \#id\\


/photos&Поиск фотоснимков вулканов\\


/tasks&Каталог результатов моделирования\\


/tasks/\#id&Информация о задаче с~идентификатором \#id\\


\hline


\end{tabular}


\end{center}


%\end{table}


%\begin{table}\small %tabl3


\begin{center}


\Caption{Параметры запроса для поиска в~архиве фотоснимков}


\vspace*{2ex}





\begin{tabular}{|l|l|}


\hline


\multicolumn{1}{|c|}{Параметр}&\multicolumn{1}{c|}{Назначение}\\


\hline


volcano\_id&Список идентификаторов вулканов, разделенных запятыми\\


\hline


camera\_id&Список идентификаторов видеокамер, разделенных запятыми\\


\hline


time&\tabcolsep=0pt\begin{tabular}{l}


Дата и~время, для которых ведется поиск,\\ задается в~формате YYYYMMDDHHIISS, 


где\\


 YYYY~--- 4~цифры года,\\ MM~--- номер


 месяца с~ведущим нулем,\\ DD~--- день месяца 


с~ведущим нулем,\\


 HH~--- часы с~ведущим нулем,\\ II~--- минуты с~ведущим нулем,\\


  SS~---  секунды с~ведущим нулем\end{tabular}\\


\hline


\end{tabular}


\end{center}


\end{table}





     При поиске фотоснимков также указываются параметры GET-запроса 


(табл.~3).


     








     Поиск фотоснимков в~архиве ведется по заданным вулканам 


и~видеокамерам для временного периода $[\mbox{time}\hm - 1{,}5~\mbox{ч}; 


\mbox{time}\hm + 1{,}5~\mbox{ч}]$.


     


При успешном выполнении запроса результат возвращается в~формате JSON. 


Он является текстовым, имеет простую (по сравнению, например, с~XML) 


разметку, поддерживается большим числом языков программирования 


и~программных библиотек. На рис.~4 представлены примеры взаимодействия 


с~сервисом.











     С применением разработанных интерфейсов реализовано взаимодействие 


между ИС VolSatView и~<<Сигнал>>. Все доступные структурированные 


наборы научных данных и~результаты работы специализированных подсистем 


ИС <<Сигнал>> (см.\ табл.~2) передаются в~ИС VolSatView, где представляются 


уже в~составе собственных пользовательских интерфейсов и~программных 


решений. Это позволяет, в~част\-ности, проводить совместный анализ данных 


результатов компьютерного моделирования траекторий распространения 


пепловых облаков и~фактических данных наблюдений, полученных средствами 


дистанционного зондирования Земли~\cite{7-kor, 8-kor}. 





\begin{figure} %fig4


 \vspace*{1pt}


 \begin{center}


 \mbox{%


 \epsfxsize=125mm 


 \epsfbox{kor-4.eps}


 }


\end{center}


 \vspace*{-11pt}


\Caption{Пример взаимодействия с~веб-сер\-ви\-сом для получения информации: 


(\textit{а})~по объектам наблюдения; (\textit{б})~по сети наблюдений; 


(\textit{в})~из архива 


изображений}


\vspace*{-6pt}


\end{figure}





\vspace*{-4pt}





\section{Заключение}





\vspace*{-2pt}





     Разработанный RESTful веб-сервис позволил сформировать 


универсальный инструмент для централизованного доступа к~разнородным 


научным данным в~области исследования и~оперативного мониторинга 


состояния вулканов Камчатки и~Северных Курил. В~рамках исследования были 


реализованы про\-грам\-мные средства для обмена данными и~работы с~ресурсами 


ведущих специализированных ИС. Это позволило 


расширить качественные и~количественные возможности прикладных систем, 


создать новые средства для совместного анализа результатов обработки 


различных видов оперативных наблюдений и~справочной информации.


     


     Дальнейшее развитие веб-сервиса связано с~подключением новых 


источников данных, содержащих информацию о результатах исследований 


вулканов Камчатки и~Северных Курил, а также с~внедрением новых средств 


обработки и~представления инструментальных данных. 





{\small\frenchspacing


{%\baselineskip=10.8pt


\addcontentsline{toc}{section}{Литература}


\begin{thebibliography}{99}


\bibitem{1-kor}


\Au{Гордеев Е.\,И., Гирина О.\,А.}


Вулканы и их опасность для авиации~/\!\!/


Вестник РАН, 2014. Т.~84. №\,2. С.~134.


 doi: 10.7868\!/\!S0869587314020121.


 





\bibitem{2-kor}


\Au{Сорокин~А.\,А., Макогонов~С.\,В., Королев~С.\,П.}


Информационная инфраструктура для коллективной работы ученых Дальнего 


Востока России~/\!\!/


Научно-тех\-ни\-че\-ская информация. Сер.~1: 


Организация и методика информационной работы, 2017.


 №\,12. С.~14--16.


 





\bibitem{3-kor}


\Au{Романова И.\,М., Гирина О.\,А., Максимов~А.\,П., Мелекесцев~И.\,В.} Создание 


комплексной информационной веб-сис\-те\-мы <<Вулканы Ку\-ри\-ло-Кам\-чат\-ской 


островной дуги>> (VOKKIA)~/\!\!/ Информатика и~системы управления, 2012. №\,3(33). 


С.~179--187.


\bibitem{4-kor}


\Au{Гирина О.\,А., Гордеев~Е.\,И.} Проект KVERT~--- снижение вулканической опас\-ности 


для авиации при эксплозивных извержениях вулканов Камчатки и~Северных Курил~/\!\!/ 


Вестник ДВО РАН, 2007. №\,2. С.~100--109.


\bibitem{5-kor}


\Au{Королёв С.\,П., Сорокин~А.\,А., Верхотуров~А.\,Л., 


Коновалов~А.\,В., Шестаков~Н.\,В.}


Автоматизированная информационная система для работы с инструментальными


 данными региональной сети сейсмологических наблюдений ДВО РАН~/\!\!/   


Сейсмические приборы, 2014.


Т.~50. №\,3. С.~30--41.


 





\bibitem{6-kor}


\Au{Sorokin A.\,A., Korolev~S.\,P., Romanova~I.\,M., Girina~O.\,A., Urmanov~I.\,P.} The 


Kamchatka volcano video monitoring system~/\!\!/ 6th Workshop (International) on Computer 


Science and Engineering.~--- Tokyo, Japan, 2016. P.~734--737.


\bibitem{7-kor}


\Au{Гордеев Е.\,И., Гирина~О.\,А., Лупян~Е.\,А., Сорокин~А.\,А., 


Крамарева~Л.\,С., Ефремов~В.\,Ю., Кашницкий~А.\,В., Уваров~И.\,А., 


Бурцев~М.\,А., Романова~И.\,М., Мельников~Д.\,В., Маневич~А.\,Г., 


Королев~С.\,П., Верхотуров~А.\,Л.}


Информационная система VOLSATVIEW для решения 


задач мониторинга вулканической активности Камчатки и~Курил~/\!\!/


 Вулканология и сейсмология, 2016.


№\,6. С.~62--77.


 doi: 10.7868\!/\!S0203030616060043.


 


\bibitem{9-kor}


\Au{Richardson L., Amundsen~M., Ruby~S.} RESTful web APIs.~--- O'Reilly Media, 2013. 408~p.





\bibitem{10-kor}


\Au{Safronov M., Winesett~J.} Web application development with Yii2 and PHP.~--- Packt 


Publishing, 2014. 406~p.











\bibitem{8-kor}


\Au{Сорокин А.\,А., Королев~С.\,П., Гирина~О.\,А., Балашов~И.\,В., Ефремов~В.\,Ю., 


Романова~И.\,М., Мальковский~С.\,И.} Интегрированная программная платформа для 


комплексного анализа распространения пепловых шлейфов при эксплозивных извержениях 


вулканов Камчатки~/\!\!/ Современные проблемы дистанционного зондирования Земли из 


космоса, 2016. Т.~13. №\,4. С.~9--19. doi:  


10.21046\!/\!2070-7401-2016-13-12-9-19.





        \end{thebibliography}


} }








\vspace*{-9pt}





\hfill{\small\textit{Поступила в~редакцию 10.03.18}}














%\vspace*{8pt}





%\hrule





%\vspace*{2pt}





\newpage





\vspace*{-28pt}





%\hrule





\vspace*{-3pt}











\def\tit{SERVICE-ORIENTED INTERFACE TO~ACCESS SCIENTIFIC DATA 


FOR~STUDY AND~STATE OPERATIONAL MONITORING 


OF~VOLCANOES OF~KAMCHATKA AND~NORTHERN KURILES\\[-5pt]}





\def\titkol{Service-oriented interface to~access scientific data 


for study and~state % operational 


monitoring} 


%of~volcanoes of~Kamchatka and~Northern Kuriles}








\def\aut{S.\,P.~Korolev$^1$, I.\,M.~Romanova$^2$, S.\,I.~Malkovsky$^1$, 


and~A.\,A.~Sorokin$^1$\\[-5pt]}


\def\autkol{S.\,P.~Korolev, I.\,M.~Romanova, S.\,I.~Malkovsky, 


and~A.\,A.~Sorokin}











\titel{\tit}{\aut}{\autkol}{\titkol}





\vspace*{-14pt}








\noindent


     $^1$Computing Center, Far Eastern Branch of the Russian Academy of Sciences,\linebreak 


$\hphantom{^1}$65~Kim~U~Chen Str., Khabarovsk 680000, Russian Federation


     


     \noindent


     $^2$Institute of Volcanology and Seismology, Far Eastern Branch of the Russian\linebreak 


$\hphantom{^1}$Academy of Sciences, 9~Piip Blvd, Petropavlovsk-Kamchatsky 683006, 


Russian\linebreak 


$\hphantom{^1}$Federation





\def\leftfootline{\small{\textbf{\thepage}


\hfill Sistemy i Sredstva Informatiki~---


Systems and Means of Informatics 2018 vol~28 no\,2}


}%


 \def\rightfootline{\small{Sistemy i Sredstva Informatiki~---


 Systems and Means of Informatics   2018   vol~28   no\,2


\hfill \textbf{\thepage}}}





 %\vspace*{-1pt}     


 





\Abste{On the territory of the Russian Far East, there are dozens of active volcanoes, 


requiring continuous attention of scientists for analysis and control of their condition. 


To provide solutions of various scientific problems, different information systems 


were implemented. However, disunity of information systems and limited data 


exchange between them limit the possibility of carrying out complex studies. This 


may result in disastrous consequences for the population and different fields of 


people's activities. The article describes the developed service-oriented software 


interface that implements interaction between the main existing information systems. 


The mutual use of accumulated sets of scientific data and processing instruments 


improves research and operational monitoring of the state of volcanoes in Kamchatka 


and Northern Kuriles.}








\KWE{volcanoes of Kamchatka; volcano monitoring; Web services;  


REST-service; data integration; databases}





\DOI{10.14357\!/\!08696527180207}





\vspace*{-16pt}





\Ack





\vspace*{-8pt}





\noindent


The work was supported by the Russian Foundation for


Basic Research (grant No.\,16-37-00026~mol\_a) 


and the Comprehensive Program of Basic Scientific Research 


of the Far Eastern Branch of the Russian Academy of Sciences ``Far East'' 


(project 18-5-091). To ensure the functioning of scientific information services, 


the resources of the Datacenter of the Far Eastern Branch of the Russian 


Academy of Sciences have been used.





 


\vspace*{-7pt}





\renewcommand{\bibname}{\protect\rmfamily References}


%\renewcommand{\bibname}{\large\protect\rm References}





{\small\frenchspacing


{%\baselineskip=10.8pt


\addcontentsline{toc}{section}{References}


\begin{thebibliography}{99}





\vspace*{-2pt}





\bibitem{1-kor-1}


\Aue{Gordeev, E.\,I., and O.\,A.~Girina.} 2014. Volcanoes and their hazard to 


aviation. \textit{Her. Rus. Acad. Sci.} 84(1):1--8. doi: 


10.1134\!/\!S1019331614010079.


\bibitem{2-kor-1}


\Aue{Sorokin, A.\,A., S.\,I.~Makogonov, and S.\,P.~Korolev.} 2017. The 


information infrastructure for collective scientific work in the Far East of Russia. 


\textit{Scientific Technical Information Processing} 44(4):302--304. doi: 


10.3103\!/\!S0147688217040153.


\bibitem{3-kor-1}


\Aue{Romanova, I.\,M., O.\,A.~Girina, A.\,P.~Maksimov, and I.\,V.~Melekestsev}. 


2012. So\-zda\-nie kompleksnoy informatsionnoy veb-sistemy ``Vulkany 


Kurilo-Kamchatskoy ost\-rov\-noy dugi'' (VOKKIA) [Creation of complex infomation 


Web system ``Volcanoes of the Kurile-Kamchatka island arc'' (VOKKIA)]. 


\textit{Informatika i~sistemy upravleniya} [Information Science and Control 


Systems] 3(33):179--187.


\bibitem{4-kor-1}


\Aue{Girina, O.\,A., and E.\,I.~Gordeev.} 2007. Proekt KVERT~--- snizhenie 


vulkanicheskoy opasnosti dlya aviatsii pri eksplozivnykh izverzheniyakh vulkanov 


Kamchatki i~Severnykh Kuril [KVERT project: Reduction of volcanic hazards for 


aviation from explosive eruptions of Kamchatka and Northern Kuriles volcanoes]. 


\textit{Vestnik of the Far East Branch of the Russian Academy of Sciences} 2:100--109.


\bibitem{5-kor-1}


\Aue{Korolev, S.\,P., A.\,A.~Sorokin, A.\,L.~Verkhoturov, A.\,V.~Konovalov, and 


N.\,V.~Shestakov.} 2015. Automated information system for instrument data 


processing of the regional seismic observation network of FEB RAS. \textit{Seism. 


Instr.} 51(3):209--218. doi: 10.3103\!/\!S0747923915030068.


\bibitem{6-kor-1}


\Aue{Sorokin, A.\,A., S.\,P.~Korolev, I.\,M.~Romanova, O.\,A.~Girina, and 


I.\,P.~Urmanov.} 2016. The Kamchatka volcano video monitoring system. \textit{6th 


Workshop (International) on Computer Science and Engineering}. Tokyo, Japan. 


734--737.


\bibitem{7-kor-1}


\Aue{Gordeev, E.\,I., O.\,A.~Girina, E.\,A.~Lupyan, A.\,A.~Sorokin, 


L.\,S.~Kramareva, V.\,Yu.~Efremov, A.\,V.~Kashnitskii, I.\,A.~Uvarov, 


M.~Burtsev, I.\,M.~Romanova, D.\,V.~Mel'nikov, A.\,G.~Manevich, 


S.\,P.~Korolev, and A.\,L.~Verkhoturov.} 2016. The VolSatView information 


system for monitoring the volcanic activity in Kamchatka and on the Kuril 


Islands. \textit{J.~Volcanol. Seismol.} 10(6):382--394. doi: 


10.1134\!/ S074204631606004X.





\bibitem{9-kor-1}


\Aue{Richardson, L., M.~Amundsen, and S.~Ruby}. 2013. \textit{RESTful web 


APIs}. O'Reilly Media. 408~p.


\bibitem{10-kor-1}


\Aue{Safronov, M., and J.~Winesett.} 2014. \textit{Web application development 


with Yii2 and PHP}. Packt Publishing. 406~p.





\bibitem{8-kor-1} %10


\Aue{Sorokin, A.\,A., S.\,P.~Korolev, O.\,A.~Girina, I.\,V.~Balashov, 


V.\,Yu.~Efremov, I.\,M.~Ro\-ma\-no\-va, and S.\,I.~Malkovsky.} 2016. Integrirovannaya 


programmnaya platforma dlya kompleksnogo analiza rasprostraneniya peplovykh 


shleyfov pri eksplozivnykh izverzheniyakh vulkanov Kamchatki [The integrated 


software platform for a comprehensive analysis of ash plume propagation from 


explosive eruptions of Kamchatka volcanoes]. \textit{Sovremennye problemy 


distantsionnogo zondirovaniya Zemli iz kosmosa} [Current Problems in Remote 


Sensing of the Earth from Space] 13(4):9--19. doi:  


10.21046\!/\!2070-7401-2016-13-12-9-19.


\end{thebibliography}


} }








\vspace*{-9pt}





\hfill{\small\textit{Received March 10, 2018}}    





\vspace*{-10pt}





\Contr





\vspace*{-4pt}





\noindent


\textbf{Korolev Sergey P.} (b.\ 1985)~--- scientist, Computing Center, Far Eastern 


Branch of the Russian Academy of Sciences, 65~Kim~U~Chen Str., Khabarovsk 


680000, Russian Federation; \mbox{serejk@febras.net}





\vspace*{1pt}





\noindent


\textbf{Romanova Iraida M.} (b.\ 1961)~--- leading programmer, Institute of 


Volcanology and Seismology, Far Eastern Branch of the Russian Academy of Sciences, 


9~Piip Blvd, Petropavlovsk-Kamchatsky 683006, Russian Federation; 


\mbox{roman@kscnet.ru}





\vspace*{1pt}





\noindent


\textbf{Malkovsky Sergei I.} (b.\ 1983)~--- scientist, Computing Center, Far Eastern 


Branch of the Russian Academy of Sciences, 65~Kim~U~Chen Str., Khabarovsk 


680000, Russian Federation; \mbox{sergey.malkovsky@gmail.com}








\vspace*{1pt}





\noindent


\textbf{Sorokin Aleksei A.} (b.\ 1980)~--- leading scientist, Computing Center, Far 


Eastern Branch of the Russian Academy of Sciences, 65~Kim~U~Chen Str., Khabarovsk 


680000, Russian Federation; \mbox{alsor@febras.net}








\label{end\stat}








\renewcommand{\bibname}{\protect\rm Литература}


